\documentclass[11pt,a4paper]{article}
% preamble.tex — shared setup for human-verification documents
% Compile each document with:  tectonic <File>.tex

\usepackage[T1]{fontenc}
\usepackage[utf8]{inputenc}
\usepackage{lmodern}
\usepackage{amssymb,amsmath,amsthm}
\usepackage{mathtools}
\usepackage[dvipsnames]{xcolor}
\usepackage{listings}
\usepackage{hyperref}
\usepackage[margin=2.5cm]{geometry}
\usepackage{parskip}

% ── Colours ──────────────────────────────────────────────────────────
\definecolor{leanBg}{HTML}{F7F7F0}
\definecolor{haskellBg}{HTML}{F0F4F7}
\definecolor{keyword}{HTML}{7B3294}
\definecolor{comment}{HTML}{5A8C32}
\definecolor{noteFg}{HTML}{B35900}

% ── Lean 4 listings language ─────────────────────────────────────────
% NOTE: Lean listing blocks use ASCII approximations of unicode symbols
% (e.g. -> for arrow, o. for compose).  The literate key substitutes
% these with proper math-mode renderings in the PDF.
\lstdefinelanguage{Lean4}{
  morekeywords={def,theorem,structure,where,namespace,end,import,fun,
                let,have,show,by,exact,intro,cases,rfl,subst,obtain,
                funext,simp,rw,apply,variable,Type,Prop,instance,
                section,open,class,extends,deriving,sorry,noncomputable},
  sensitive=true,
  morecomment=[l]{--},
  morecomment=[n]{/-}{-/},
  morestring=[b]",
}

% ── Lean listing style ───────────────────────────────────────────────
\lstdefinestyle{lean}{
  language=Lean4,
  backgroundcolor=\color{leanBg},
  basicstyle=\small\ttfamily,
  keywordstyle=\color{keyword}\bfseries,
  commentstyle=\color{comment}\itshape,
  breaklines=true,
  frame=single,
  framerule=0.4pt,
  xleftmargin=4pt,
  xrightmargin=4pt,
  aboveskip=6pt,
  belowskip=4pt,
  literate=
    {->}{{$\to$}}2
    {=>}{{$\Rightarrow$}}2
    {<|}{{$\langle$}}1
    {|>}{{$\rangle$}}1
    {==>}{{$\Longrightarrow$}}3
    { o. }{{$\;\circ\;$}}3
    {forall}{{$\forall$}}1
    {:=}{{$\,:=\,$}}2
    {!=}{{$\neq$}}2
    {<=}{{$\leq$}}2
    {>=}{{$\geq$}}2,
}

% ── Haskell listing style ───────────────────────────────────────────
\lstdefinestyle{haskell}{
  language=Haskell,
  backgroundcolor=\color{haskellBg},
  basicstyle=\small\ttfamily,
  keywordstyle=\color{keyword}\bfseries,
  commentstyle=\color{comment}\itshape,
  breaklines=true,
  frame=single,
  framerule=0.4pt,
  xleftmargin=4pt,
  xrightmargin=4pt,
  aboveskip=6pt,
  belowskip=4pt,
  literate=
    {->}{{$\to$}}2
    {=>}{{$\Rightarrow$}}2,
}

% ── Theorem environments ────────────────────────────────────────────
\theoremstyle{definition}
\newtheorem{defn}{Definition}[section]
\newtheorem{thm}[defn]{Theorem}
\newtheorem{prop}[defn]{Proposition}

\theoremstyle{remark}
\newtheorem*{fidelitynote}{Fidelity note}

% ── Convenience macros ───────────────────────────────────────────────
\newcommand{\lean}[1]{\texttt{#1}}      % inline Lean identifier
\newcommand{\hask}[1]{\texttt{#1}}      % inline Haskell identifier
\newcommand{\Ext}{\mathrm{Ext}}
\newcommand{\Pos}{\mathrm{Pos}}
\newcommand{\Shape}{\mathrm{Shape}}


% Double-bracket notation for the extension functor
\usepackage{stmaryrd}

\title{Human Verification: \texttt{Containers.Dirichlet}\\[4pt]
  \large The Dirichlet Tensor $(\mathbf{Cont}, \otimes, y)$}
\author{Machine-generated from \texttt{lean/Dirichlet.lean} for review by Neil Ghani}
\date{July 2026}

\begin{document}
\maketitle

\begin{abstract}
This document presents every definition and theorem \emph{statement} from the
Lean~4 module \lean{Containers.Dirichlet} in three parallel forms: the original
Lean, a readable Haskell rendering, and precise mathematical English.  Proofs
are omitted --- they are machine-checked by Lean.  The purpose is to let a
human verifier confirm that the formalisation captures the intended
mathematics.

The file defines the \textbf{Dirichlet tensor} $C \otimes D$ (shapes $S \times T$,
positions $P\,s \times Q\,t$), proves it is a monoidal structure on $\mathbf{Cont}$
(all coherences by \lean{rfl}), and discusses when $\llbracket{-}\rrbracket$ is
(or is not) strong monoidal for $\otimes$.  A uniformity morphism $C \otimes D
\to C \triangleleft D$ is defined, and shown to be a proper mono (not an iso).
\end{abstract}

\tableofcontents
\bigskip

%===================================================================
\section{Overview}
%===================================================================

For containers $C = S \triangleleft P$ and $D = T \triangleleft Q$, the Dirichlet
tensor is:
\[
  C \otimes D = (S \times T) \triangleleft (\lambda\,(s,t).\; P\,s \times Q\,t).
\]
Contrast the categorical product:
\[
  C \times D = (S \times T) \triangleleft (\lambda\,(s,t).\; P\,s \sqcup Q\,t).
\]
The shapes are the same; only the positions differ --- product vs.\ sum.

The key asymmetry is in extensions:
\begin{align*}
  \llbracket C \times D \rrbracket X
    &= \textstyle\Sigma_{(s,t)},\; X^{P\,s \sqcup Q\,t}
     \cong \llbracket C \rrbracket X \times \llbracket D \rrbracket X
    &\text{(pointwise, by } y^{A+B} \cong y^A \times y^B), \\
  \llbracket C \otimes D \rrbracket X
    &= \textstyle\Sigma_{(s,t)},\; X^{P\,s \times Q\,t}
     = \textstyle\Sigma_{(s,t)},\; (X^{P\,s})^{Q\,t}
    &\text{(not pointwise; } y^{A \times B} \text{ does not split)}.
\end{align*}

%===================================================================
\section{The tensor and its unit}
%===================================================================

\subsection{Definition: \lean{Container.dirichlet}}

\begin{defn}[Dirichlet tensor]
For containers $C$ and $D$, the \textbf{Dirichlet tensor} $C \otimes D$ has:
\begin{itemize}
  \item shapes: $S_C \times S_D$,
  \item positions at $(s,t)$: $P_C(s) \times P_D(t)$.
\end{itemize}
\end{defn}

\paragraph{Lean.}
\begin{lstlisting}[style=lean]
def Container.dirichlet (C D : Container) : Container where
  Shape := C.Shape * D.Shape
  Pos   := fun st => C.Pos st.1 * D.Pos st.2
\end{lstlisting}

\paragraph{Haskell.}
\begin{lstlisting}[style=haskell]
dirichlet :: Container s p -> Container t q -> Container (s, t) (p, q)
-- Shape = (s, t);  Pos (s, t) = (p, q)
\end{lstlisting}


\subsection{Definition: \lean{Container.y}}

\begin{defn}[Unit for $\otimes$]
The unit for the Dirichlet tensor is
$y = \mathbf{1} \triangleleft (\lambda\,\ast.\,\mathbf{1})$
--- the same container as $I$, the unit for $\triangleleft$.  In the Dirichlet
setting it is read as the \emph{representable} $y^1$.
\[
  y \otimes C \text{ has shapes } \mathbf{1} \times S \cong S
  \text{ and positions } \mathbf{1} \times P\,s \cong P\,s.
\]
\end{defn}

\paragraph{Lean.}
\begin{lstlisting}[style=lean]
abbrev Container.y : Container := Container.I
\end{lstlisting}

\paragraph{Haskell.}
\begin{lstlisting}[style=haskell]
-- y = I = the container ((), const ())
-- Unit for both otimes and triangleleft.
containerY :: Container () ()
containerY = Container () (const ())
\end{lstlisting}


%===================================================================
\section{Bifunctoriality}
%===================================================================

\subsection{Definition: \lean{Container.dir\textsubscript{2}}}

\begin{defn}[Bifunctorial action of $\otimes$]
For morphisms $\mu : C \Rightarrow C'$ and $\nu : D \Rightarrow D'$, the action
$\lean{dir}_2\,\mu\,\nu : C \otimes D \Rightarrow C' \otimes D'$ is
\emph{componentwise}:
\begin{itemize}
  \item on shapes: $(s,t) \mapsto (\mu(s), \nu(t))$,
  \item on positions (backward): $(s,t),\,(p,q) \mapsto (\mu^\sharp_s(p), \nu^\sharp_t(q))$.
\end{itemize}
\end{defn}

\paragraph{Lean.}
\begin{lstlisting}[style=lean]
def Container.dir2 {C C' D D' : Container}
    (mu : ContainerMorphism C C') (nu : ContainerMorphism D D') :
    ContainerMorphism (C |*| D) (C' |*| D') where
  onShapes := fun st => (mu.onShapes st.1, nu.onShapes st.2)
  onPos    := fun st pq => (mu.onPos st.1 pq.1, nu.onPos st.2 pq.2)
\end{lstlisting}

\paragraph{Haskell.}
\begin{lstlisting}[style=haskell]
dir2 :: Morph c p c' p' -> Morph d q d' q'
     -> Morph (c `Dir` d) (p, q) (c' `Dir` d') (p', q')
dir2 mu nu = Morph
  { onShapes = \(s, t) -> (onShapes mu s, onShapes nu t)
  , onPos    = \(s, t) (p, q) -> (onPos mu s p, onPos nu t q) }
\end{lstlisting}

\begin{fidelitynote}
Unlike $\triangleleft$, the Dirichlet morphism action has no outer-into-inner
threading: both coordinates are treated independently.  This makes the functor
laws hold by \lean{rfl} --- a purely definitional computation.
\end{fidelitynote}


\subsection{Theorem: \lean{Container.dir\textsubscript{2}\_id}}

\begin{thm}[Dirichlet bifunctoriality: identity]
$\lean{dir}_2\,(\mathrm{id}_C)\,(\mathrm{id}_D) = \mathrm{id}_{C \otimes D}$.
Holds by \lean{rfl}.
\end{thm}

\paragraph{Lean.}
\begin{lstlisting}[style=lean]
theorem Container.dir2_id (C D : Container) :
    Container.dir2 (ContainerMorphism.id C) (ContainerMorphism.id D)
      = ContainerMorphism.id (C |*| D) := rfl
\end{lstlisting}


\subsection{Theorem: \lean{Container.dir\textsubscript{2}\_comp}}

\begin{thm}[Dirichlet bifunctoriality: interchange]
$\lean{dir}_2\,(\mu \cdot \mu')\,(\nu \cdot \nu')
= \lean{dir}_2\,\mu\,\nu \;\cdot\; \lean{dir}_2\,\mu'\,\nu'$.
Holds by \lean{rfl}.
\end{thm}

\paragraph{Lean.}
\begin{lstlisting}[style=lean]
theorem Container.dir2_comp {C C' C'' D D' D'' : Container}
    (mu : ContainerMorphism C C') (mu' : ContainerMorphism C' C'')
    (nu : ContainerMorphism D D') (nu' : ContainerMorphism D' D'') :
    Container.dir2 (mu.comp mu') (nu.comp nu')
      = (Container.dir2 mu nu).comp (Container.dir2 mu' nu') := rfl
\end{lstlisting}


%===================================================================
\section{Structural isomorphisms}
%===================================================================

\subsection{Definition: \lean{Container.dirAssociator}}

\begin{defn}[Dirichlet associator]
The isomorphism $(C \otimes D) \otimes E \cong C \otimes (D \otimes E)$:
\begin{itemize}
  \item Shapes: $\mathrm{Prod}$-reassociation $((s,t),u) \mapsto (s,(t,u))$.
  \item Positions (backward): also $\mathrm{Prod}$-reassociation, running backwards:
    given $(p,(q,r))$, produce $((p,q),r)$ (and vice versa).
\end{itemize}
Both round trips hold by \lean{rfl} ($\eta$ for $\mathrm{Prod}$).
\end{defn}

\paragraph{Lean.}
\begin{lstlisting}[style=lean]
def Container.dirAssociator (C D E : Container) :
    ContainerIso ((C |*| D) |*| E) (C |*| (D |*| E)) where
  hom :=
    { onShapes := fun x => (x.1.1, (x.1.2, x.2))
      onPos    := fun _ p => ((p.1, p.2.1), p.2.2) }
  inv :=
    { onShapes := fun x => ((x.1, x.2.1), x.2.2)
      onPos    := fun _ p => (p.1.1, (p.1.2, p.2)) }
  hom_inv := rfl
  inv_hom := rfl
\end{lstlisting}

\paragraph{Haskell.}
\begin{lstlisting}[style=haskell]
dirAssociator ::
  ContainerIso (Dir (Dir c d) e) (Dir c (Dir d e))
-- Shapes: ((s,t),u) <-> (s,(t,u))
-- Positions (backward): ((p,q),r) <-> (p,(q,r))
\end{lstlisting}


\subsection{Definition: \lean{Container.dirLeftUnitor}}

\begin{defn}[Dirichlet left unitor]
The isomorphism $y \otimes C \cong C$:
\begin{itemize}
  \item Forward on shapes: $(\ast, s) \mapsto s$.
  \item Forward on positions (backward): $p \mapsto (\ast, p)$.
  \item Backward on shapes: $s \mapsto (\ast, s)$.
  \item Backward on positions (backward): $(\ast, p) \mapsto p$.
\end{itemize}
Both round trips hold by \lean{rfl}.
\end{defn}

\paragraph{Lean.}
\begin{lstlisting}[style=lean]
def Container.dirLeftUnitor (C : Container) :
    ContainerIso (Container.y |*| C) C where
  hom :=
    { onShapes := fun x => x.2
      onPos    := fun _ p => ((), p) }
  inv :=
    { onShapes := fun s => ((), s)
      onPos    := fun _ p => p.2 }
  hom_inv := rfl
  inv_hom := rfl
\end{lstlisting}

\paragraph{Haskell.}
\begin{lstlisting}[style=haskell]
dirLeftUnitor :: ContainerIso (Dir YContainer c) c
-- fwd shapes: ((), s) -> s
-- bwd positions: p -> ((), p)
\end{lstlisting}


\subsection{Definition: \lean{Container.dirRightUnitor}}

\begin{defn}[Dirichlet right unitor]
The isomorphism $C \otimes y \cong C$:
\begin{itemize}
  \item Forward on shapes: $(s, \ast) \mapsto s$.
  \item Forward on positions (backward): $p \mapsto (p, \ast)$.
\end{itemize}
Both round trips hold by \lean{rfl}.
\end{defn}

\paragraph{Lean.}
\begin{lstlisting}[style=lean]
def Container.dirRightUnitor (C : Container) :
    ContainerIso (C |*| Container.y) C where
  hom :=
    { onShapes := fun x => x.1
      onPos    := fun _ p => (p, ()) }
  inv :=
    { onShapes := fun s => (s, ())
      onPos    := fun _ p => p.1 }
  hom_inv := rfl
  inv_hom := rfl
\end{lstlisting}


%===================================================================
\section{Naturality and coherence}
%===================================================================

\subsection{Theorem: \lean{Container.dirAssociator\_naturality}}

\begin{thm}[Naturality of Dirichlet associator]
\[
  (\lean{dir}_2\,(\lean{dir}_2\,f_1\,f_2)\,f_3) \cdot \alpha^\otimes_{Y_1,Y_2,Y_3}
  \;=\;
  \alpha^\otimes_{X_1,X_2,X_3} \cdot \lean{dir}_2\,f_1\,(\lean{dir}_2\,f_2\,f_3).
\]
Holds by \lean{rfl}.
\end{thm}

\paragraph{Lean.}
\begin{lstlisting}[style=lean]
theorem Container.dirAssociator_naturality
    {X1 Y1 X2 Y2 X3 Y3 : Container}
    (f1 : ContainerMorphism X1 Y1)
    (f2 : ContainerMorphism X2 Y2)
    (f3 : ContainerMorphism X3 Y3) :
    (Container.dir2 (Container.dir2 f1 f2) f3).comp
        (Container.dirAssociator Y1 Y2 Y3).hom
      = (Container.dirAssociator X1 X2 X3).hom.comp
          (Container.dir2 f1 (Container.dir2 f2 f3)) := rfl
\end{lstlisting}


\subsection{Theorem: \lean{Container.dirLeftUnitor\_naturality}}

\begin{thm}[Naturality of Dirichlet left unitor]
$(\lean{dir}_2\,(\mathrm{id}_y)\,f) \cdot \lambda^\otimes_Y = \lambda^\otimes_X \cdot f$.
Holds by \lean{rfl}.
\end{thm}

\paragraph{Lean.}
\begin{lstlisting}[style=lean]
theorem Container.dirLeftUnitor_naturality {X Y : Container}
    (f : ContainerMorphism X Y) :
    (Container.dir2 (ContainerMorphism.id Container.y) f).comp
        (Container.dirLeftUnitor Y).hom
      = (Container.dirLeftUnitor X).hom.comp f := rfl
\end{lstlisting}


\subsection{Theorem: \lean{Container.dirRightUnitor\_naturality}}

\begin{thm}[Naturality of Dirichlet right unitor]
$(\lean{dir}_2\,f\,(\mathrm{id}_y)) \cdot \rho^\otimes_Y = \rho^\otimes_X \cdot f$.
Holds by \lean{rfl}.
\end{thm}

\paragraph{Lean.}
\begin{lstlisting}[style=lean]
theorem Container.dirRightUnitor_naturality {X Y : Container}
    (f : ContainerMorphism X Y) :
    (Container.dir2 f (ContainerMorphism.id Container.y)).comp
        (Container.dirRightUnitor Y).hom
      = (Container.dirRightUnitor X).hom.comp f := rfl
\end{lstlisting}


\subsection{Theorem: \lean{Container.dirPentagon}}

\begin{thm}[Pentagon for $\otimes$]
The two re-association paths of
$((C_1 \otimes C_2) \otimes C_3) \otimes C_4 \to
  C_1 \otimes (C_2 \otimes (C_3 \otimes C_4))$ agree.
Holds by \lean{rfl}.
\end{thm}

\paragraph{Lean.}
\begin{lstlisting}[style=lean]
theorem Container.dirPentagon (C1 C2 C3 C4 : Container) :
    ((Container.dir2 (Container.dirAssociator C1 C2 C3).hom
          (ContainerMorphism.id C4)).comp
        (Container.dirAssociator C1 (C2 |*| C3) C4).hom).comp
        (Container.dir2 (ContainerMorphism.id C1)
          (Container.dirAssociator C2 C3 C4).hom)
      = (Container.dirAssociator (C1 |*| C2) C3 C4).hom.comp
          (Container.dirAssociator C1 C2 (C3 |*| C4)).hom := rfl
\end{lstlisting}


\subsection{Theorem: \lean{Container.dirTriangle}}

\begin{thm}[Triangle for $\otimes$]
$\alpha^\otimes_{C_1,y,C_2} \cdot (\mathrm{id}_{C_1} \otimes \lambda^\otimes_{C_2})
= \rho^\otimes_{C_1} \otimes \mathrm{id}_{C_2}$.
Holds by \lean{rfl}.
\end{thm}

\paragraph{Lean.}
\begin{lstlisting}[style=lean]
theorem Container.dirTriangle (C1 C2 : Container) :
    (Container.dirAssociator C1 Container.y C2).hom.comp
        (Container.dir2 (ContainerMorphism.id C1)
          (Container.dirLeftUnitor C2).hom)
      = Container.dir2 (Container.dirRightUnitor C1).hom
          (ContainerMorphism.id C2) := rfl
\end{lstlisting}


\subsection{Definition: \lean{contDirichletMonoidal}}

\begin{defn}[$(\mathbf{Cont}, \otimes, y)$ is a monoidal category]
A named \lean{def}: the Dirichlet tensor $\otimes$ with morphism action
$\lean{dir}_2$, unit $y$, and the structural isomorphisms above satisfy all
monoidal-category axioms.  Every law holds by \lean{rfl}.
\end{defn}

\paragraph{Lean.}
\begin{lstlisting}[style=lean]
@[reducible] def contDirichletMonoidal : MonoidalCategory Container where
  tensorObj         := Container.dirichlet
  tensorHom         := Container.dir2
  tensorHom_id      := Container.dir2_id
  tensorHom_comp    := fun f1 g1 f2 g2 => Container.dir2_comp f1 g1 f2 g2
  tensorUnit        := Container.y
  associator        := fun X Y Z => (Container.dirAssociator X Y Z).toCatIso
  leftUnitor        := fun X => (Container.dirLeftUnitor X).toCatIso
  rightUnitor       := fun X => (Container.dirRightUnitor X).toCatIso
  associator_naturality  :=
    fun f1 f2 f3 => Container.dirAssociator_naturality f1 f2 f3
  leftUnitor_naturality  := fun f => Container.dirLeftUnitor_naturality f
  rightUnitor_naturality := fun f => Container.dirRightUnitor_naturality f
  pentagon          := Container.dirPentagon
  triangle          := Container.dirTriangle
\end{lstlisting}

\paragraph{Haskell.}
\begin{lstlisting}[style=haskell]
contDirichletMonoidal :: MonoidalCategoryDef ContainerMorph
-- tensor = dirichlet, unit = y, all laws by rfl.
\end{lstlisting}


%===================================================================
\section{Extension and the Dirichlet tensor}
%===================================================================

\subsection{Theorem: \lean{Container.dirichletExt\_normalForm}}

\begin{thm}[Normal form for the Dirichlet extension]
By definition (i.e.\ by \lean{rfl}):
\[
  \llbracket C \otimes D \rrbracket X
  \;=\;
  \Sigma_{(s,t) : S \times T},\; (P\,s \times Q\,t \to X).
\]
This is a \emph{computation / unfolding rule}, not a strong-monoidality statement.
The right-hand side mentions the \emph{presentation} $S, T, P, Q$, not the
functors $\llbracket C \rrbracket$ and $\llbracket D \rrbracket$.  The correct
intrinsic account is that $\llbracket{-}\rrbracket$ is \emph{strong} (not
strict) monoidal for $\otimes$, the target tensor being the Day convolution of
$({\bf Set}, \times, \mathbf{1})$ (Niu--Spivak, Prop.~3.79).
\end{thm}

\paragraph{Lean.}
\begin{lstlisting}[style=lean]
theorem Container.dirichletExt_normalForm
    (C D : Container) (X : Type) :
    Ext (C |*| D) X
      = ((st : C.Shape * D.Shape) * (C.Pos st.1 * D.Pos st.2 -> X)) := rfl
\end{lstlisting}

\paragraph{Haskell.}
\begin{lstlisting}[style=haskell]
-- By definition:
-- Ext (dirichlet c d) x  =  ((s,t), (p,q) -> x)
-- This is NOT the same as Ext c x `tensor` Ext d x for any pointwise tensor.
\end{lstlisting}

\begin{fidelitynote}
Three of the four monoidal structures give genuine (non-identity) comparison
isomorphisms for $\llbracket{-}\rrbracket$: $\llbracket G \triangleleft F
\rrbracket X \cong \llbracket G \rrbracket (\llbracket F \rrbracket X)$,
$\llbracket C + D \rrbracket X \cong \llbracket C \rrbracket X \sqcup
\llbracket D \rrbracket X$, and $\llbracket C \times D \rrbracket X \cong
\llbracket C \rrbracket X \times \llbracket D \rrbracket X$.  For $\otimes$
there is no such pointwise comparison, witnessed by $y^3 \otimes y^3 = y^9$
while any pointwise tensor would give $y^6$.
\end{fidelitynote}


%===================================================================
\section{The uniformity morphism \texorpdfstring{$C \otimes D \to C \triangleleft D$}{C tensor D -> C seq D}}
%===================================================================

\subsection{Definition: \lean{Container.dirToSeq}}

\begin{defn}[Uniformity morphism]
The canonical map $C \otimes D \Rightarrow C \triangleleft D$ sends the shape
$(s, t)$ to $(s, \lambda\,{-}.\;t)$ (the constant inner-shape assignment):
\[
  (s, t) \;\longmapsto\; (s,\; \lambda p.\; t).
\]
On positions: the identity, since the fibres agree definitionally:
$\Sigma\,p : P\,s,\; Q((\lambda\,{-}.\;t)\,p) = P\,s \times Q\,t$.
\end{defn}

\paragraph{Lean.}
\begin{lstlisting}[style=lean]
def Container.dirToSeq (C D : Container) :
    ContainerMorphism (C |*| D) (C seq D) where
  onShapes := fun st => <|st.1, fun _ => st.2|>
  onPos    := fun _ pq => (pq.1, pq.2)
\end{lstlisting}

\paragraph{Haskell.}
\begin{lstlisting}[style=haskell]
dirToSeq :: ContainerMorphism (c `Dir` d) (c `Seq` d)
dirToSeq = Morph
  { onShapes = \(s, t) -> (s, const t)  -- constant inner assignment
  , onPos    = \_ (p, q) -> (p, q) }
\end{lstlisting}


\subsection{Helper definitions: \lean{Container.ySq}, \lean{Container.two}, \lean{Container.seqShapeSnd}}

\begin{defn}[Auxiliary containers for the counterexample]
\begin{itemize}
  \item $y^2 = \mathbf{1} \triangleleft (\lambda\,\ast.\,\mathrm{Bool})$:
    one shape, two positions.
  \item $2 = \mathrm{Bool} \triangleleft (\lambda\,{-}.\,\emptyset)$:
    two shapes, no positions.
  \item \lean{seqShapeSnd}: extracts the inner-shape function from a shape of
    $y^2 \triangleleft 2$, viewed as a plain function $\mathrm{Bool} \to
    \mathrm{Bool}$.  Defined to enable \lean{congrArg} in the counterexample.
\end{itemize}
\end{defn}

\paragraph{Lean.}
\begin{lstlisting}[style=lean]
def Container.ySq : Container where
  Shape := Unit
  Pos   := fun _ => Bool

def Container.two : Container where
  Shape := Bool
  Pos   := fun _ => Empty

def Container.seqShapeSnd
    (z : (Container.ySq seq Container.two).Shape) : Bool -> Bool := z.2
\end{lstlisting}

\paragraph{Haskell.}
\begin{lstlisting}[style=haskell]
ySq :: Container () Bool   -- one shape, two positions
ySq = Container { shape = (), pos = const False {- placeholder -} }

two :: Container Bool Void  -- two shapes, no positions
two = Container { shapes = [True, False], pos = const absurd }
\end{lstlisting}


\subsection{Theorem: \lean{Container.dirToSeq\_not\_surjective}}

\begin{thm}[The uniformity morphism is not an isomorphism]
The shape map of \lean{dirToSeq} for $C = y^2$, $D = 2$ is
$\mathrm{Bool} \ni t \mapsto (\ast, \lambda\,{-}.\;t)$, which hits only
\emph{constant} functions $\mathrm{Bool} \to \mathrm{Bool}$.  The identity
function $\mathrm{id} : \mathrm{Bool} \to \mathrm{Bool}$ is not in the image.
Therefore \lean{dirToSeq} is not surjective on shapes, hence not an
isomorphism.

Concretely, $y^2 \otimes 2$ has 2 shapes while $y^2 \triangleleft 2$ has 4,
so the extensions differ: $\llbracket y^2 \otimes 2 \rrbracket X$ has 2
elements and $\llbracket y^2 \triangleleft 2 \rrbracket X$ has 4.
Thus $\otimes \neq \triangleleft$.
\end{thm}

\paragraph{Lean.}
\begin{lstlisting}[style=lean]
theorem Container.dirToSeq_not_surjective :
    not (Exists fun st : (Container.ySq |*| Container.two).Shape =>
        (Container.dirToSeq Container.ySq Container.two).onShapes st
          = (<|(), fun b => b|> : (Container.ySq seq Container.two).Shape))
\end{lstlisting}

\paragraph{Haskell.}
\begin{lstlisting}[style=haskell]
-- The image of dirToSeq on shapes is { ((), const t) | t :: Bool }.
-- The shape ((), id) of (ySq `Seq` two) is not in this image,
-- since id True = True /= False = id False but const t True = const t False.
-- So dirToSeq is not surjective, hence not an iso.
\end{lstlisting}

\begin{fidelitynote}
The proof works by contradiction: assuming a preimage $st$ exists,
$\lean{congrArg}\,\lean{seqShapeSnd}$ extracts the equation
$\lambda\,{-}.\;st.2 = \mathrm{id}$ on $\mathrm{Bool}$; applying at
$\mathrm{true}$ gives $st.2 = \mathrm{true}$ and at $\mathrm{false}$ gives
$st.2 = \mathrm{false}$, a contradiction.
\end{fidelitynote}

\end{document}
